\documentclass[11pt]{article}
% ===========================================================================
%  Shared preamble for the AP Calculus AB exam-preparation worksheets.
%  Input right after \documentclass. Files beginning with "_" are not compiled
%  on their own. Navy and gold, to match the course's lesson pages.
% ===========================================================================
\usepackage[letterpaper,top=2.3cm,bottom=2.2cm,left=1.7cm,right=1.7cm,headheight=28pt,footskip=30pt]{geometry}
\usepackage{amsmath,amssymb}
\usepackage[T1]{fontenc}
\usepackage[utf8]{inputenc}
\usepackage{lmodern}
\usepackage{xcolor}
\usepackage[most]{tcolorbox}
\usepackage{enumitem}
\usepackage{fancyhdr}
\usepackage{array}
\usepackage{tikz}
\usepackage{pgfplots}
\pgfplotsset{compat=1.18}

\definecolor{navy}{HTML}{161E54}
\definecolor{ink}{HTML}{27357E}
\definecolor{gold}{HTML}{B98A00}
\definecolor{goldbg}{HTML}{FFFDF3}
\definecolor{inkbg}{HTML}{F6F7FD}
\definecolor{rule}{HTML}{C9CFF0}
\definecolor{okgreen}{HTML}{1B7A44}
\definecolor{okbg}{HTML}{F1F8F3}

\pagestyle{fancy}
\fancyhf{}
\fancyhead[L]{\footnotesize Name: \underline{\hspace{4cm}}}
\fancyhead[C]{\textbf{\color{navy}Integration Academy}\\[-2pt]{\scriptsize Calculus AB \textperiodcentered{} AP\textsuperscript{\textregistered} Exam Preparation}}
\fancyfoot[L]{\tiny\color{black!55}\parbox[t]{0.78\textwidth}{AP\textsuperscript{\textregistered} is a registered trademark of the College Board, which was not involved in the production of, and does not endorse, this material. \textcopyright{} Integration Academy.}}
\fancyfoot[R]{\scriptsize\thepage}
\renewcommand{\headrulewidth}{1.2pt}
\renewcommand{\headrule}{\hbox to\headwidth{\color{gold}\leaders\hrule height \headrulewidth\hfill}}
\renewcommand{\footrulewidth}{0pt}
\setlength{\parindent}{0pt}
\setlength{\parskip}{4pt}

% Title block: \wstitle{code}{title}{unit line}
\newcommand{\wstitle}[3]{\begin{center}{\small\bfseries\color{gold}WORKSHEET #1}\\[3pt]{\LARGE\bfseries\color{navy}#2}\\[3pt]{\small\color{black!60}#3}\end{center}\vspace{2pt}}
% Section heading
\newcommand{\sech}[1]{\par\vspace{10pt}{\large\bfseries\color{navy}#1}\par\nobreak\vspace{2pt}{\color{rule}\hrule height0.8pt}\vspace{7pt}}
% A level tag printed beside a group of questions
\newcommand{\level}[1]{\par\vspace{6pt}{\footnotesize\bfseries\color{gold}\MakeUppercase{#1}}\par\nobreak\vspace{2pt}}

\newtcolorbox{ideabox}{colframe=ink,colback=inkbg,coltitle=white,colbacktitle=ink,fonttitle=\bfseries\small,title={Key ideas},arc=2pt,boxrule=0.8pt,left=8pt,right=8pt}
\newtcolorbox{example}[1]{colframe=rule,colback=white,coltitle=navy,colbacktitle=inkbg,fonttitle=\bfseries\small,title={#1},arc=2pt,boxrule=0.8pt,borderline west={3pt}{0pt}{ink},left=10pt,right=8pt}
\newtcolorbox{frbox}[1]{colframe=gold,colback=goldbg,coltitle=white,colbacktitle=gold,fonttitle=\bfseries\small,title={#1},arc=2pt,boxrule=0.8pt,breakable,left=8pt,right=8pt}
\newtcolorbox{solbox}{colframe=okgreen,colback=okbg,coltitle=white,colbacktitle=okgreen,fonttitle=\bfseries\small,title={Worked solutions},arc=2pt,boxrule=0.8pt,breakable,left=8pt,right=8pt}
\newtcolorbox{scorebox}[1]{colframe=navy,colback=white,coltitle=white,colbacktitle=navy,fonttitle=\bfseries\small,title={#1},arc=2pt,boxrule=0.8pt,breakable,left=8pt,right=8pt}

\newcommand{\soln}{\par\smallskip\textbf{\color{ink}Solution.}\ }
\newcommand{\result}[1]{\par\smallskip\textbf{#1}}
% blank working space of a given height
\newcommand{\work}[1]{\par\vspace{#1}}
% a point value in the margin of a scoring guide
\newcommand{\pts}[1]{\hfill{\footnotesize\bfseries\color{gold}[#1]}}

% Coordinate grid for graph questions: \begin{apgraph}{xmin}{xmax}{ymin}{ymax} ... \end{apgraph}
\newenvironment{apgraph}[4]{\begin{center}\begin{tikzpicture}\begin{axis}[
  axis lines=middle, xmin=#1, xmax=#2, ymin=#3, ymax=#4,
  grid=both, grid style={black!12}, xtick distance=1, ytick distance=1,
  ticklabel style={font=\scriptsize}, xlabel={$x$}, ylabel={$y$},
  width=10cm, height=7.2cm, samples=80, clip=false,
  every axis plot/.append style={line width=1.1pt, color=ink}]}{\end{axis}\end{tikzpicture}\end{center}}
\newcommand{\closeddot}[2]{\addplot[only marks, mark=*, mark size=2.2pt, color=ink] coordinates {(#1,#2)};}
\newcommand{\opendot}[2]{\addplot[only marks, mark=*, mark size=2.2pt, color=ink, mark options={fill=white}] coordinates {(#1,#2)};}

% Multiple-choice options: four across, or stacked when the options are long.
\newcommand{\opts}[4]{\par\vspace{3pt}\begin{tabular}{@{}p{0.235\linewidth}p{0.235\linewidth}p{0.235\linewidth}p{0.235\linewidth}@{}}(A)\ #1 & (B)\ #2 & (C)\ #3 & (D)\ #4\end{tabular}\par\vspace{9pt}}
\newcommand{\optsv}[4]{\par\vspace{3pt}\begin{tabular}{@{}l@{\ \ }p{0.9\linewidth}@{}}(A) & #1\\(B) & #2\\(C) & #3\\(D) & #4\end{tabular}\par\vspace{9pt}}

\fancyhead[R]{\footnotesize\color{black!60}Worksheet 1.9}
\begin{document}

\wstitle{1.9}{Unbounded Behaviour: Vertical and Horizontal Asymptotes}{Unit 1: Limits and Continuity}

\begin{ideabox}
\begin{itemize}[leftmargin=*,itemsep=2pt,topsep=2pt]
  \item $x=c$ is a \textbf{vertical asymptote} if a one-sided limit at $c$ is $\pm\infty$. Find the sign factor by factor.
  \item $y=L$ is a \textbf{horizontal asymptote} if $f(x)\to L$ as $x\to\infty$ or $x\to-\infty$.
  \item Rational functions at infinity: compare degrees; equal degrees give the ratio of leading coefficients.
  \item $\sqrt{x^2}=|x|$, which equals $-x$ when $x\to-\infty$. Growth: $\ln x\ll x^p\ll e^x$.
\end{itemize}
\end{ideabox}

\sech{Worked examples}
\begin{example}{Example 1 \textperiodcentered{} The sign of an infinite limit}
Find $\displaystyle\lim_{x\to 1^-}\dfrac{x+3}{x-1}$ and $\displaystyle\lim_{x\to 1^+}\dfrac{x+3}{x-1}$.\soln The numerator is near 4. For $x<1$ the denominator is a small negative number; for $x>1$ a small positive number.\result{$-\infty$ from the left; $\infty$ from the right. The line $x=1$ is a vertical asymptote.}
\end{example}

\begin{example}{Example 2 \textperiodcentered{} Rational limits at infinity}
Evaluate $\displaystyle\lim_{x\to \infty}\dfrac{2x^2-7}{3x^2+x}$ and $\displaystyle\lim_{x\to -\infty}\dfrac{x+1}{x^2+4}$.\soln First: equal degrees, ratio of leading coefficients. Second: the denominator has the higher degree.\result{$\dfrac23$ and $0$.}
\end{example}

\begin{example}{Example 3 \textperiodcentered{} A square root as $x\to-\infty$}
Evaluate $\displaystyle\lim_{x\to -\infty}\dfrac{\sqrt{x^2+9}}{2x-1}$.\soln $\sqrt{x^2+9}=|x|\sqrt{1+\tfrac{9}{x^2}}=-x\sqrt{1+\tfrac9{x^2}}$ for $x<0$. The expression is $\dfrac{-\sqrt{1+\frac9{x^2}}}{2-\frac1x}$.\result{The limit is $-\dfrac12$.}
\end{example}

\newpage
\sech{Practice}

\level{Core}
\begin{enumerate}[leftmargin=*,itemsep=2pt]
  \item Evaluate $\displaystyle\lim_{x\to \infty}\dfrac{5x^3-x}{2x^3+4}$.\work{2.3cm}
  \item Evaluate $\displaystyle\lim_{x\to -2^+}\dfrac{1}{x+2}$.\work{2.3cm}
  \item Evaluate $\displaystyle\lim_{x\to \infty}\dfrac{4x+1}{x^2-3}$.\work{2.3cm}
  \item Evaluate $\displaystyle\lim_{x\to -\infty}e^x$.\work{2.3cm}
\end{enumerate}

\level{Exam level}
\begin{enumerate}[leftmargin=*,itemsep=2pt,start=5]
  \item Find both horizontal asymptotes of $f(x)=\dfrac{4e^x-1}{2e^x+3}$.\work{2.8cm}
  \item Find all vertical and horizontal asymptotes, and any holes, of $f(x)=\dfrac{x^2-4}{x^2-5x+6}$.\work{2.8cm}
  \item Evaluate $\displaystyle\lim_{x\to -\infty}\dfrac{3x+2}{\sqrt{4x^2+1}}$.\work{2.8cm}
  \item Evaluate $\displaystyle\lim_{x\to \infty}\left(\sqrt{x^2+8x}-x\right)$.\work{2.8cm}
\end{enumerate}

\begin{frbox}{Question 9 \textperiodcentered{} Free-response practice}
Let $f(x)=\dfrac{3x^2-12}{x^2+x-6}$.
\begin{enumerate}[label=(\alph*),leftmargin=*,itemsep=2pt]
  \item Show that $f$ has a removable discontinuity at $x=2$, and find $\displaystyle\lim_{x\to 2}f(x)$.\work{2.6cm}
  \item Write an equation for the vertical asymptote of the graph of $f$. Justify with a limit.\work{2.6cm}
  \item Write an equation for the horizontal asymptote of the graph of $f$. Justify with a limit.\work{2.6cm}
\end{enumerate}
\end{frbox}

\newpage
\begin{solbox}
\begin{enumerate}[leftmargin=*,itemsep=4pt]
  \item $\dfrac52$.
  \item $\infty$.
  \item $0$.
  \item $0$.
  \item As $x\to\infty$: $\dfrac{4-e^{-x}}{2+3e^{-x}}\to2$. As $x\to-\infty$: $\dfrac{0-1}{0+3}=-\dfrac13$. Asymptotes $y=2$ and $y=-\dfrac13$.
  \item $f(x)=\dfrac{(x-2)(x+2)}{(x-2)(x-3)}=\dfrac{x+2}{x-3}$ for $x\neq2$. Vertical asymptote $x=3$; horizontal asymptote $y=1$; hole at $(2,-4)$.
  \item $\sqrt{4x^2+1}=-x\sqrt{4+\tfrac1{x^2}}$ for $x<0$, so the limit is $\dfrac{3}{-2}=-\dfrac32$.
  \item $\dfrac{8x}{\sqrt{x^2+8x}+x}=\dfrac{8}{\sqrt{1+\frac8x}+1}\to4$.
  \item \begin{enumerate}[label=(\alph*),leftmargin=*,itemsep=3pt]
    \item $f(x)=\dfrac{3(x-2)(x+2)}{(x+3)(x-2)}=\dfrac{3(x+2)}{x+3}$ for $x\neq2$. So $\displaystyle\lim_{x\to 2}f(x)=\dfrac{12}{5}$ exists while $f(2)$ is undefined: removable.
    \item $x=-3$, because $\displaystyle\lim_{x\to -3^+}f(x)=-\infty$ (numerator near $-3$, denominator small and positive).
    \item $y=3$, because $\displaystyle\lim_{x\to \infty}f(x)=3$.
  \end{enumerate}
\end{enumerate}
\end{solbox}

\end{document}
